\documentclass[12pt]{article}

%% Basic document formatting
\usepackage{amsmath, amsthm, amssymb, amsfonts}
\usepackage{mathrsfs}
\usepackage{mathtools}
\usepackage{xspace}
\usepackage{thmtools}
\usepackage{graphicx}
\usepackage{tikz}
\usepackage{tikz-cd} 
\usepackage{ytableau}
\usepackage{blkarray}
\usepackage{hyperref}
\usepackage{caption}
\usepackage{setspace}
\usepackage{array}
\usepackage{fancyhdr}
\usepackage{titling}
\usepackage[left=0.4in,right=0.4in,top=1in,bottom=1in]{geometry}
\usepackage{float}
\usepackage{cancel}
\usepackage{tabularx}
\usepackage[utf8]{inputenc}
\usepackage[english]{babel}
\usepackage{framed}
\usepackage[dvipsnames]{xcolor}
\usepackage{environ}
\usepackage{tcolorbox}
\tcbuselibrary{theorems,skins,breakable}

\usepackage{enumitem}
\setlist[enumerate]{leftmargin=*}
\setlist[enumerate,1]{labelindent=\parindent}
\setlist[enumerate,2]{labelindent=0pt}

% Blackboard Bold
\newcommand{\Z}{\mathbb{Z}}                 % integers
\newcommand{\Zpl}{\mathbb{Z}_{+}}           % positive integers
\newcommand{\Znn}{\mathbb{Z}_{\geq 0}}      % nonnegative integers. 
\newcommand{\Q}{\mathbb{Q}}                 % rationals
\newcommand{\Qpl}{\mathbb{Q}_{+}}           % positive Rationals
\newcommand{\R}{\mathbb{R}}                 % reals
\newcommand{\Rpl}{\mathbb{R}_{+}}           % positive Reals
\newcommand{\C}{\mathbb{C}}                 % complex numbers
\newcommand{\F}{\mathbb{F}}                 % field

% Words
\newcommand{\st}{\text{ such that }}
\newcommand{\wrt}{\text{ with respect to }}
\newcommand{\with}{\text{ with }}
\newcommand{\ie}{\text{, i.e. }}
\newcommand*{\ditto}{\text{---}\texttt{"}\text{---}}

% Operators
\newcommand{\abs}[1]{\left|#1\right|}                   % absolute value
\newcommand{\norm}[1]{\abs{\abs{#1}}}
\newcommand{\floor}[1]{\left\lfloor #1 \right\rfloor}   % floor
\newcommand{\ceil}[1]{\left\lceil #1 \right\rceil}      % ceiling
\newcommand{\powset}[1]{\mathscr{P}(#1)}

% Category Theory 
\renewcommand{\hom}{\text{Hom}}
\newcommand{\homspace}[3]{\hom_{#1}(#2, #3)}
\newcommand{\coproduct}[9]{
  \begin{tikzcd}[ampersand replacement=\&]
      \& \& #1 \arrow[dll, bend right, "#6"] \arrow[dl, "#5"] \\
   #4 \& #3 \arrow[l, "#9"] \\
      \& \& #2 \arrow[ull, bend left, "#8"] \arrow[ul, "#7"]
  \end{tikzcd}
}

\newcommand{\product}[9]{ 
  \begin{tikzcd}[ampersand replacement=\&]
      \& \& #3 \\
      #1 \arrow[r, "#7"] \arrow[urr, bend left, "#5"] \arrow[drr, bend right, "#6"] \& #2 \arrow[ur, "#8"] \arrow[dr, "#9"] \\ 
      \& \& #4
  \end{tikzcd}
}


% Algebra 
\newcommand{\Syl}[2]{\text{Syl}_{#1}{(#2)}}           % set of Sylow-p subgroups 
\newcommand{\<}{\langle}                % \<x\>, subgroup generated by x 
\renewcommand{\>}{\rangle}
\renewcommand{\index}[2]{[#1 : #2]}
\newcommand{\id}{\text{id}}             % identity element
\newcommand{\lhdneq}{%
  \mathrel{\ooalign{$\lneq$\cr\raise.22ex\hbox{$\lhd$}\cr}}}
\newcommand{\order}[1]{\text{o}(#1)}    % order of an element
\newcommand{\spec}{\operatorname{Spec}}
\newcommand{\rad}{\operatorname{rad}}   % radical of an ideal 
\newcommand{\sym}[1]{\mathfrak{S}_{#1}}
\newcommand{\aut}[1]{\text{Aut}(#1)} 
\newcommand{\gal}[2]{\text{Gal}(#1 / #2)} 
\newcommand{\inn}[1]{\text{Inn}(#1)} 
\let\oldcong\cong
\let\oldequiv\equiv
\renewcommand{\cong}{\oldequiv}
\renewcommand{\equiv}{\oldcong}

\newcommand{\triangleleftneq}{\mathrel{\ooalign{%
  $\trianglelefteq$\cr
  \hidewidth\raise0.06ex\hbox{$\not\mkern4mu$}\hidewidth\cr
}}}

\makeatletter % cycle
\newcommand{\cyc}[1]{(\mathbf{\cyc@process#1\relax})}
\def\cyc@process#1#2\relax{%
	#1%
	\ifx\relax#2\relax
	\else
		\,\cyc@process#2\relax
	\fi
}
\makeatother

\newcommand{\GL}[2]{\text{GL}_{#1}(#2)}                             % general linear group
\newcommand{\SL}[2]{\text{SL}_{#1}(#2)}                             % special linear group
\newcommand{\gl}[2]{\mathfrak{gl}_{#1}(#2)}
\renewcommand{\sl}[2]{\mathfrak{sl}_{#1}(#2)}
\newcommand{\so}[2]{\mathfrak{so}_{#1}(#2)}
\newcommand{\su}[1]{\mathfrak{su}(#1)}

% Linear Algebra
\newcommand{\bmat}[1]{\begin{bmatrix}#1\end{bmatrix}}   % bracketed matrix
\newcommand{\rank}{\operatorname{rank}}                 % rank
\newcommand{\nullity}{\operatorname{nullity}}           % nullity
\newcommand{\tr}{\operatorname{tr}}

% Combinatorics
\newcommand{\qnum}[1]{\left[#1\right]_q}                % q-analog
\newcommand{\qfact}[1]{\left[#1\right]!_q}              % q-analog factorial 
\newcommand{\qbinom}[2]{\binom{#1}{#2}_q}               % Gaussian binomial coefficient

% Topology/Analysis
\newcommand{\ball}[2]{\text{B}_{#1}(#2)}  % B_r(x): open r-balls around x
\newcommand{\diam}{\text{diam}}           % diamter of a set in metric space 
\newcommand{\lowint}[2]{\underline{\int_{#1}^{#2}}}
\newcommand{\upint}[2]{\overline{\int}_{#1}^{#2}}

% Misc Notation
\newcommand{\oneton}[1]{\{1, \ldots, #1\}}
\newcommand{\defeq}{\vcentcolon=}   % :=
\newcommand{\eqdef}{=\vcentcolon}   % =: 
\renewcommand{\bf}[1]{\textbf{#1}}
\newcommand{\im}{\operatorname{im}}
\newcommand{\fnrest}[2]{\left.#1\right|_{#2}}
\newcommand{\isom}{\overset{\sim}{\rightarrow}} 


\renewcommand{\thesection}{\S \arabic{section}}
\renewcommand{\thesubsection}{\arabic{subsection}.}
\renewcommand{\thesubsubsection}{(\alph{subsubsection})} 

\newtheorem{claim}{Claim}
\newtheorem*{lemma}{Lemma}

%% Headers & title setup
\newcommand{\course}{\textit{Algebra: Chapter 0}, Aluffi} 
\newcommand{\myname}{Jay Ser}
\setlength{\headheight}{14.5pt}
\pagestyle{fancy}
\fancyhf{}
\renewcommand{\headrulewidth}{0.4pt}
\lhead{\course}
\rhead{\myname}
\cfoot{\thepage}
\setlength{\droptitle}{-4em} 
\title{\course\ \\ Part III: Rings and Modules}
\author{\myname}
\date{April 18, 2026 $\sim$ \today}

\begin{document}
\maketitle
\thispagestyle{fancy}

\newcommand{\dirsum}[2]{#1^{\oplus#2}}

%------------------------------------------------------------------------------%

Notation: 
\begin{itemize}
  \item Given a ring $R$ and element $a \in R$, I denote by $\<a\>$ the principal \textit{ideal} generated by $a$ in the ring $R$ and $\<a\>^{+}$ the \textit{subgroup} generated by $a$ in the underlying additive group $R^{+}$. 
\end{itemize}

\setcounter{section}{2}
\section{Ideals and Quotient Rings}
\setcounter{subsection}{3} 
\subsection{}
Let $\varphi: \Z \rightarrow R$ be the unique ring homomorphism from $\Z$ to $R$. 
By definition of a ring homomorphism, $\varphi(1_{\Z}) = 1_{R}$. 
By assumption on $R$, the subgroup $\<1_{R}\>^{+}$ of the additive group $R^{+}$ is equal to the ideal $\<1_R\>$ of $R$. 
$\<1_{R}\> = R$, of course. 
Also, notice $\im \varphi = \<1_{R}\>^{+}$. 
Hence $\im \varphi = R$. 

By definition, $n \defeq \text{char}R$ is the kernal of $\varphi: \Z \rightarrow R$, hence the first isomorphism theorem yields 
$$\Z / n\Z \equiv R.$$

\subsection{} 
Let $e_{ij} \in \mathcal{M}_{n}(R)$ denote the matrix with a 1 on the $(i, j)$th entry and 0's everywhere else. 
Suppose $A \in J$. 
Since $J$ is a two-sided ideal, 
$$e_{rs}Ae_{kl} \in J, \hspace{1.5em} (\forall r, s, k, l \in \oneton{n})$$
What is this the matrix $e_{rs} A e_{kl}$? 
Let $A = [a_{ij}]$. 
Writing $[a'_{ij}] \defeq Ae_{kl}$, 
\begin{align*}
  a'_{ij} & = \sum_{u = 1}^{n} a_{iu} (e_{kl})_{uj} \\ 
          & = \begin{cases} 
            0, & \text{if } j \neq l \\ 
            a_{ik}, & \text{if } j = l.
          \end{cases}
\end{align*}
Next, writing $[a''_{ij}] = e_{rs}[a'_{ij}]$, 
\begin{align*}
  a''_{ij} & = \sum_{u = 1}^{n} (e_{rs})_{iu} a'_{uj} \\ 
           & = \begin{cases} 
            0, & \text{if } j \neq l \text{ or } i \neq r \\ 
            a'_{sj} = a_{sk}, & \text{if } j = l \text{ and } i = r
           \end{cases}
\end{align*}
This shows that $e_{rs} A e_{kl}$ is the matrix whose $(r, l)$th entry is the $(s, k)$th entry of $A$ and has 0's everywhere else.
This proves the forward direction. 

Next, suppose every entry of $a$ of $A$, there is a matrix $a e_{ij} \in J$ for all $i, j \in \oneton{n}$. 
Then clearly $A \in J$: 
writing $A = [a_{ij}]$, 
$$A = \sum_{1 \leq i, j \leq n} a_{ij} e_{ij}.$$
and because each $a_{ij} e_{ij} \in J$, so is their sum, 
i.e., $A \in J$. 

\setcounter{subsection}{11}
\subsection{} 
The nilradical of a commutative ring is an ideal because the binomial theorem holds in commutative rings. 

$\mathcal{M}_{n}(\R)$ for $n > 1$ is a noncommutative ring whose set of nilpotent elements is not an ideal.
Recall that lower/upper triangular matrices with 0's on the diagonal are nilpotent. 
However, the sums of such matrices may be an invertible matrix; 
the powers of an invertible matrix are of course cannot be the zero matrix. 
For example, in $\mathcal{2}{\R}$, $\bmat{0 & 1 \\0 & 0}$ and $\bmat{0 & 0 \\1 & 0}$ are nilpotent. 
However, their sum is $\bmat{0 & 1 \\ 1 & 0}$, which is an invertible matrix. 
This demonstrates that the set of nilpotent elements of $\mathcal{M}_{2}(\R)$ does not form an ideal. 

\setcounter{subsection}{13}
\subsection{}
Suppose that $\text{char} R = nm$ for integers $n, m > 1$. 
In $R$, let $n \defeq \underbrace{1 + \ldots + 1}_{n \text{ times}}$ and $m \defeq \underbrace{1 + \ldots + 1}_{m \text{ times}}$. 
Then by the assumption on the character of $R$, 
$$n, m \neq 0, \text{ but } nm = 0,$$
which violates the fact that $R$ is an integral domain. 
Hence the character of an integral doamin is either 0 or a prime integer. 

Clearly, the zero ring is the only ring with character 1. 

%TODO: ba \neq ab doesn't necessarily lead to contradiction when R is not an integral domain!!
\subsection{} 
$1 + 1 = (1 + 1)(1 + 1) = 1 + 1 + 1 + 1$ shows that $1 + 1 = 0$ since addition can be cancelled in any ring;
the character of a boolean ring is 2. 
Next, for all $a, b \in R$, 
\begin{align*}
  a + b & = (a + b)^2 \\ 
        & = a^2 + ab + ba + b^2 \\ 
        & = a + b + ab + ba,
\end{align*}
so $ab + ba = 0$, 
i.e., $ab = -ba$. 
Because $R$ has character 2, $ba = -ba$, from which follows $ab = ba$;
$R$ is commutative. 

Suppose any boolean ring is an integral domain. 
Then for any $a \neq 0$, the cancellation law applies to the identity $a^2 = a$ to show that $a = 1$. 
Hence any boolean ring that is an integral domain is isomorphic to $\Z / 2\Z$. 

%TODO: do Q4.17 after learning some topology, moron. 
\section{Ideals and Quotients: Remarks and Examples}
\setcounter{subsection}{8}
\subsection{} 
Since $f(x)$ is a zerodivisor, choose $g(x) \in R[x]$ of minimal degree such that $fg = 0$. 
Write $f(x) = a_n x^n + \ldots + a_0$ and $g(x) = b_m x^m + \ldots + b_0$ with coefficients in $R$.  
Also, let $f(x)g(x) = c_{n + m} x^{n + m} + \ldots + 0$, where each $c_{t} = 0$ by assumption. 
I claim that each $a_i$ for $i \in \{0, \ldots, n\}$ annihilates $g$, that is, 
$$a_i g(x) = 0 \hspace{1.5em} (i \in \{0, \ldots, n\})$$
Induct on the $i$, starting from $i = n$. 
Using the regular rules for polynomial multiplication, 
$$c_{n + m} = a_n b_m,$$
and $c_{n + m} = 0$ by assumption. 
Hence $a_n b_m = 0$. 
Because $R$ is commutative, $f (a_n g(x)) = a_n f(x) g(x) = 0$, and $a_n b_m = 0$ means $a_{n} g(x)$ is a polynomial of lower degree than $g$.  
The minimality of $g$ thus forces 
$$a_n g(x) = 0;$$
$a_n$ annihilates $g$. 

Next, assume $a_{n}, \ldots, a_{n - k + 1}$ annihilate $g$. 
Once again, starting from the assumption that $c_{n + m - k} = 0$ and using the regular rules for polynomial multiplication, 
\begin{align*}
  0 = c_{n + m - k} & = \sum_{i = n-k}^{n} a_{i}b_{n + m - k - i} \\ 
                    & = a_{n - k}b_m,
\end{align*}
where the last inequality follows from the inductive hypothesis. 
As in the base case, $a_{n - k}g(x)$ is a polynomial of lower degree than $g(x)$ such that $f(x) (a_{n - k})g(x) = 0$, so one concludes $a_{n - k}g(x) = 0$. 

Having shown that each coefficient $a_{i}$ annihilates $g$, it is clear that $b_j f(x) = 0$ for each $j \in \{0, \ldots, m\}$. 
Namely, $b_0 f(x) = 0$, as was to be shown (this also shows that $g(x) = b_0$, i.e., a constant polynoimal). 

\setcounter{subsection}{10}
\subsection{} 
\subsubsection{}
Define the substitution map 
$$\varphi:R[x] \rightarrow R; \hspace{0.5em} \varphi_{R} = \id, \, x \mapsto a.$$
By Prop 4.6 in the book, $\ker \varphi = \<x - a\>$ and $\varphi$ is surjective. 
I now freely use a fact that technically isn't covered yet in the book: 
since $x - a$ is an irreducible polynomial in $R[x]$, 
$$R[x] / \<x - a\> \equiv R$$ 
as rings, not just as abelian groups. 
$I \defeq \<f_1(x), \ldots, f_r(x), x - a\>$ contains $\ker \varphi$, so $J \defeq \varphi(I)$ is an ideal in $R / \<x - a\>$. 
Namely, $J$ is generated by the images of each of the generators of $I$ under $\varphi$; 
$$J = \<\overline{f_1(a)}, \ldots, \overline{f_r(a)}\>.$$ 
By the correspondence theorem, $\varphi^{-1}(J) = I$, but one also sees $\varphi^{-1}(J) = \<f_1(a), \ldots, f_r(a), x - a\>$ just by taking preimage of the generators of $J$ under $\varphi$.  
It follows that $\<f_1(x), \ldots, f_r(x), x - a\> = \<f_1(a), \ldots, f_r(a), x - a\>$. 

\subsubsection{}
Letting bars denote images under the substitution homomorphism $\varphi$ as defined above, 
\begin{align*}
  \frac{R[x]}{\<f_1(x), \ldots, f_r(x), x - a\>} & \equiv \frac{R[x] / \<x - a\>}{\<\overline{f_1(x)}, \ldots, \overline{f_r(x)}, \overline{x - a}\>} \\ 
                                                 & \equiv \frac{R}{\<f_1(a), \ldots, f_r(b)\>}
\end{align*}
where the last equivalence, particularly in the denominator, follows because the image of $f_i(x)$ under $\varphi$ is just substitution by $x = a$. 

\subsection{}
Induct on $n$. 
$n = 1$ has already been proven; 
$R[x] / \<x - a\> \equiv R$ as rings. 
Suppose the proposition is true for $n - 1$. 
Consider the ring homomorphism defined by 
$$\varphi: R[x_1, \ldots, x_n] \rightarrow R[x_1, \ldots, x_{n-1}]; \hspace{0.5em} \varphi_{R[x_1, \ldots, x_{n - 1}]} = \id, \, x_n \mapsto a_n.$$
Viewing $R[x_1, \ldots, x_n]$ as $R[x_1, \ldots, x_{n - 1}][x_n]$, i.e., polynomials with coefficients in $R[x_1, \ldots, x_{n - 1}]$ with indeterminant $x_n$, example 4.7 gives $\ker \varphi = \<x_n - a_n\>$.
By the first isomorphism thoerem and prop 3.11 in the book,
\begin{align*}
  \frac{R[x_1, \ldots, x_n]}{\<x_1 - a_n, \ldots, x_n - a_n\>} 
  & \equiv \frac{R[x_1, \ldots, x_n] / \<x_n - a_n\>}{\<\overline{x_1 - a_1}, \ldots, \overline{x_{n - 1} - a_{n - 1}}, \overline{x_{n} - a_n}\>} \\ 
  & \equiv \frac{R[x_1, \ldots, x_{n - 1}]}{\<x_1 - a_1, \ldots, x_{n - 1} - a_{n - 1}\>}
\end{align*}
By the inductive hypothesis, this last quotient ring is isomorphic to $R$. 

\subsection{} 
Take any $k \in \oneton{n}$. 
By Question 12, 
$$\frac{R[x_1, \ldots, x_n]}{\<x_1, \ldots, x_k\>} \equiv R[x_{k + 1}, \ldots, x_n]$$
where the latter is an integral domain because $R$ is an integral domain. 
Hence $\<x_1, \ldots, x_{k}\>$ is a prime ideal. 

\subsection{} 
Take maximal ideal $M \lhd R$ and $ab \in M$. 
Suppose $a \notin M$. 
Since $M$ is maximal, $M + \<a\> = R$, hence $1 = m + ad$ for some $m \in M$ and $d \in R$. 
Then
$$b = b(m + ad) = mb + abd \in M,$$
as was to be shown. 

\section{Modules Over a Ring} 
\newcommand{\cat}{R\textsf{-Mod}}
\setcounter{subsection}{8}
\subsection{} 
Define 
\begin{align*}
  \Phi:R \rightarrow \text{End}_{\cat}(M); \hspace{0.5em} & r \mapsto \sigma_r \\ 
  \sigma_{r}:M \rightarrow M; \hspace{0.5em} & m \mapsto rm
\end{align*}
$\sigma_{r}$ is clearly an $R$-module homomorphism. 
Further, for all $r, s, a \in R$ and all $m \in M$,  
\begin{align*}
  \Phi(r + as)(m) & = \sigma_{r + as}(m) \\ 
                  & =  (r + as)(m) = rm + asm \\ 
               & = \sigma_{r}(m) + \sigma_{a}\sigma_{s}(m) = (\Phi(r) + \Phi(a)\Phi(s))(m),
\end{align*}
hence $\Phi$ is a ring homomorphism. 
It is clear that $\im \Phi \subset Z(\text{End}_{\cat}(M))$. 
This shows that $\text{End}_{\cat}(M)$ is an $R$-algebra. 
$(\text{End}_{\cat}(M), \Phi)$ is a natural $R$-algebra in the sense that $\Phi$ induces the following action of $R$ on $\text{End}_{\cat}(M)$:
$$\rho: R \times \text{End}_{\cat}(M) \rightarrow \text{End}_{\cat}(M); \hspace{0.5em} (r, \varphi) \mapsto r\varphi,$$
where $r \varphi:M \rightarrow M$ is the $R$-module homomorphism defined with respect to $r$ and $\varphi$ as usual. 

$\mathcal{M}_{n}(R)$ is also an $R$-algebra, given by the homomorphism 
$$\Phi: R \rightarrow \mathcal{M}_{n}(R); \hspace{0.5em} r \mapsto \bmat{r \\ & r \\ & & \ddots \\ & & & r},$$
i.e., realize $r \in R$ as the $n \times n$ scalar matrix. 
Scalar matrices clearly commute with all other matrices in $\mathcal{M}_{n}(R)$. 

\setcounter{subsection}{10} 
\subsection{}
First, take a $R$-module homomorphism $\varphi:M \rightarrow M$. 
Then define an $R[x]$-module structure on $M$ as follows:
given $f(x) = a_n x^n + \ldots + a_0 \in R[x]$, 
$$\rho: R[x] \times M \rightarrow M; \hspace{0.5em} (f(x), m) \mapsto a_n \varphi^{n}(m) + a_{n - 1} \varphi^{n - 1}(m) + \ldots + a_0 m,$$ 
i.e., the $R[x]$-module structure of $M$ is fully determined by the underlying $R$-module structure of $M$ and the mapping $\rho(x, m) = \varphi(m)$ since, by definition, $R[x]$ is generated by $R$ and $x$. 
Instead of expliciting writing $\rho$, I just use multiplication notation for the following verificaiton that this does indeed define an $R[x]$-module structure on $M$. 
It is clear that for any $f(x), g(x) \in R[x]$ and any $m_1, m_2 \in M$, $(f(x) + g(x))(m_1) = f(x)m_1 + g(x)m_1$ and $f(x)(m_1 + m_2) = f(x)m_1 + f(x)m_2$. 
Write $f(x) = a_n x^n + \ldots + a_0$ and $g(x) = b_m x^m + \ldots + b_0$. 
Without loss of generality, assume $n = m$. 
\begin{align*}
  (f(x)g(x))m & = \left(\sum_{1 \leq i, j \leq n} a_i b_j x^{i + j}\right) m \\ 
              & = \sum_{1 \leq i, j \leq n}a_i b_j \varphi^{i + j}(m) \\ 
              & = \sum_{i = 1}^{n} a_i \varphi^{i} \left(\sum_{j = 1}^{n} b_j \varphi^{j}(m)\right) \\ 
              & = f(x) (g(x)m)
\end{align*}
where the third equality follows from the fact that $\varphi$ is a $R$-module homomorphism.
This shows that $\rho$ is defines a $R[x]$-module structure on $M$. 

Vice versa, take any action $\rho: R[x] \times M \rightarrow M$ that defines an $R[x]$-module structure on $M$. 
Then one can easily check that the map
$$\varphi: M \rightarrow M; \hspace{0.5em} m \mapsto \rho(x, m)$$
is an $R$-module homomorphism. 
This establishes the desired bijection. 

\setcounter{subsection}{12} 
\subsection{}
Let $I = \<a\>$ as an ideal in $R$ for some $a \in R \setminus 0$. 
$$\varphi: R \rightarrow I; \hspace{0.5em} r \mapsto ra$$ 
is an $R$-module homomorphism: 
for all $r_1, r_2, c \in R$, $\varphi(r_1 + cr_2) = (r_1 + cr_2)a = r_1a + cr_2a = \varphi(r_1) + c\varphi(r_2)$. 
By definition, $\varphi$ is surjective. 
It is also injective because $R$ is an integral domain: 
since $a \neq 0$, $ra = 0 \iff r = 0$. 
This shows that $I \equiv R$ as $R$-modules.

\setcounter{subsection}{14}
\subsection{} 
This is a confusing problem, even though, in the end, there's not much to it. 
Consider 
$$\varphi: I \rightarrow I \cdot (R / J); \hspace{0.5em} a \mapsto a + J.$$
First, notice that this definition does indeed make sense because $I \cdot (R / J) \lhd R / J$ as a submodule. 
For any $a, b \in I$ and any $r \in R$, the computation $\varphi(a + rb) = (a + rb) + J = (a + J) + r(b + J) = \varphi(a) + r\varphi(b)$ shows that $\varphi$ is an $R$-module homomorphism.
$\varphi$ is surjective because elements of $I \cdot (R / J)$ are of the form 
$$\sum a_i (r_i + J) = \sum (a_i' + J) \hspace{1.5em} (a_i, a_i' \in I, \, r_i \in R)$$
where $a_i' = a_i r_i \in I$ since $I$ is an ideal in $R$. 
It is clear that $\ker \varphi = I \cap J$, hence 
$$I / (I \cap J) \equiv I \cdot (R / J)$$ 
as $R$-modules. 
By the second isomorphism theorem for modules, $(I + J) / J \equiv I / (I \cap J)$, so $(I + J) / J \equiv I \cdot (R / J)$.  

\subsection{} 
Clearly, if $M = 0$, then $aM = 0$ as well, hence $aM = M$. 
Conversely, suppose $aM = M$. 
Let $a^n = 0$. 
For any $m \in M$, $\exists m_1 \in M$ such that $m = am_1$. 
But for this $m_1$, $\exists m_2 \in M \st m_1 = am_2$, i.e., $m = a^2 m_2$. 
Repeating this argument, there is an $m_n \in M \st m = a^n m_n$. 
But $a^n = 0$, hence $m = 0$. 
This shows that $M = 0$. 

\section{Products, Coproducts, etc. in R-Mod}
\setcounter{subsection}{2} 
\subsection{} 
The difficulty with this question is knowing how to map $M \rightarrow \ker p$. 
Thankfully, one has inspirations from projections in linear algebra that maps a vector to its ``projection" into a fixed subspace.  
Given any vector $v$, one can use this projection map, call it $p$, to determine the component of $v$ that is orthogonal to the subspace: 
it is given by $v - p(v)$. 
Of course, projection of vectors orthogonal to the subspace is just 0, i.e., the kernal of the projection map is precisely vectors orthogonal to the subspace.  

Something similar can be done here: 
\begin{equation} \label{6_3_orthog_proj}
  o: M \rightarrow \ker p; \hspace{0.5em} m \mapsto m - p(m)
\end{equation}
gives a well-defined mapping, and in fact it is an $R$-module homomorphism because $p$ is so. 
\eqref{6_3_orthog_proj} and 
\begin{equation} \label{6_3_proj}
  p': M \mapsto \im p; \hspace{0.5em} m \mapsto p(m)
\end{equation}
can be used to demonstrate the desired $R$-module isomorphism $M \equiv \ker p \oplus \im p$. 
The direct way to show this is checking 
$$\varphi: M \mapsto \ker p \oplus \im p; \hspace{0.5em} m \mapsto (o(m), p(m))$$
is an $R$-module isomorphism; 
the only notable point in this procedure is checking the surjectivity of $\varphi$: 
for any $(n, p(m)) \in \ker p \oplus \im p$, $\varphi(p(m) + n) = (n, p(m))$. 

Alternatively, verify that $M$ with the mappings \eqref{6_3_orthog_proj} and \eqref{6_3_proj} satisfies the universal property for the product of $\im p$ and $\ker p$: 
indeed, in the category of $R$-modules, suppose the diagram below commutes for some $\sigma: K \rightarrow M$: 
\begin{center}
  \product{K}{M}{\im p}{\ker p}{f}{g}{\sigma}{p'}{o}
\end{center}
The commutativity of the diagram forces $f(k) = p'(\sigma(k))$ and $g(k) = \sigma(k) - p(\sigma(k))$ for all $k \in K$. 
Combining these two equations yields $\sigma(k) = f(k) + g(k)$ for all $k \in K$, thus giving a unique definition for $\sigma$. 
And indeed, with this definition, $\sigma: K \rightarrow M$ is an $R$-module homomorphism because $f$ and $g$ are so. 
This verifies $M \equiv \ker p \oplus \im p$. 

\newcommand{\rsum}{R^{\oplus n}}
\setcounter{subsection}{5}
\subsection{}
This question has taught me to stop relying on category theory for intuition. 
The key point in this question is finding $n$ $R$-mod homomorphisms $\rsum \rightarrow R$ that span $\homspace{R}{\rsum}{R}$. 
We know exactly the set of $n$ ``natural" morphisms $\rsum \rightarrow R$: 
the coordinate projections 
$$\pi_{i}: \rsum \rightarrow R; \hspace{0.5em} \sum_{i = 1}^{n} r_i e_i \mapsto r_i,$$
where the elements of $\rsum$ are written $\sum_{i = 1}^{n} r_i e_i$ for $e_i: [n] \rightarrow R$, $e_i(k) = \begin{cases} 1, & \text{ if } i = k \\ 0, & \text{ else} \end{cases}$. 
And indeed, they span the homspace:  
take any $\varphi \in \homspace{R}{\rsum}{R}$. 
The identity 
$$\varphi(\sum_{i = 1}^{n} r_i e_i) = \sum_{i = 1}^{n} r_i \varphi(e_i),$$
following from the $R$-linearity of $\varphi$, shows 
$$\varphi = \sum_{i = 1}^{n} \varphi(e_i) \pi_{i}.$$

Using this observation, one can easily verify that the map 
$$\iota: [n] \rightarrow \homspace{R}{\rsum}{R}; \hspace{0.5em} i \mapsto \pi_{i}$$ 
satisfies the universal property for the free group on $[n]$.
The conclusion is that $\homspace{R}{\rsum}{R} \equiv \rsum$ as $R$-modules. 
Always think abouut algebraic intuition before thinking about categories!!!

\subsection{} 
Although $A$ of course need not be a finite index set, for the sake of thinking of good categorical definitions of $\prod_{\alpha \in A} M_{\alpha}$ and $\oplus_{\alpha \in A} M_{\alpha}$, it helps to draw the respective commuting diagrams as if they were direct products/coproducts of finitely many modules.  
Unfortunately, I'm unsure how to typeset such a diagram concisely. 
Anyhow, 
\begin{align*}
  \prod_{\alpha \in A} M_\alpha \defeq \{f:A \rightarrow \bigcup_{\alpha \in A} M_\alpha & \mid \forall \alpha \in A, \, f(\alpha) \in M_\alpha\} \\ 
  \bigoplus_{\alpha \in A} M_\alpha \defeq \{f: A \rightarrow \bigcup_{\alpha \in A} M_\alpha \mid \forall \alpha \in A, \, f(\alpha) \in M_\alpha, & \text{ and } f(\alpha) \neq 0 \text{ for finitely many } \alpha \in A\}
\end{align*}
along with the standard projections for the product and inclusions for the direct sum satisfy the categorical universal property for product and coproduct, respectively. 

Given this definition, it is clear $\Z^{\oplus \Zpl} \not \equiv \Z^{\Zpl}$.
There is clearly the inclusion homomorphism $\iota: \Z^{\oplus \Zpl} \hookrightarrow \Z^{\Zpl}$, but it is also clear that this map isn't surjective. 

\setcounter{subsection}{8}
\subsection{} 
If $F$ is the free module on $A$, its elements can be written as $\sum_{a \in A} r_a e_a$, where $e_a: A \rightarrow R$ is defined $e_a(a') = \begin{cases} 1, & \text{ if } a' = a \\ 0, & \text{ else} \end{cases}$. 
In other words, $\{e_a \mid a \in A\}$ form a basis for the $R$-module $F$, and it follows from the universal property of free modules that any $R$-module homomorphism $F \rightarrow L$ is determined entirely by how the basis elements $e_a$ are mapped. 

First, assume $\varphi: M \rightarrow N$ is surjective and take any $\alpha: F \rightarrow N$. 
For each $a \in A$, denote $n_a \defeq \alpha(e_a)$. 
Because $\varphi$ is surjective, there is $m_a \in M \st \varphi(m_a) = n_a$ for each $a \in A$. 
Letting $\beta: F \rightarrow M$ be the map uniquely determined by the mappings $e_a \mapsto m_a$, one sees $\varphi \beta = \alpha$. 

Suppose that the converse holds. 
Take any $n \in N$. 
Let $\alpha: F \rightarrow N$ be an $R$-module homomorphism such that $e_a \mapsto n$ for some $a \in A$
(the existence of such a homomorphism is guaranteed by the universal property of free modules). 
Since $\exists \beta: F \rightarrow M$ such that $\alpha = \varphi \beta$, $\varphi(\beta(e_a)) = n$.
This holds for any $n \in N$, so $\varphi$ is surjective. 

\setcounter{subsection}{12} 
\subsection{}
Let $A \subset M$ generate $M$: 
there is a surjective map $\iota: \dirsum{R}{A} \twoheadrightarrow M$. 
Take any $R$-module homomorphism $\varphi: M \rightarrow N$. 
Intuitively, $\varphi(A)$ generates $\im \varphi$. 
Indeed, there are surjective mappings 
$$\dirsum{R}{A} \overset{\iota}{\twoheadrightarrow} M \overset{\varphi}{\twoheadrightarrow} \im \varphi.$$ 
Given any $\alpha = \sum_{a \in A} r_a a \in \dirsum{R}{A}$, the fact that $\iota$ and $\varphi$ are $R$-module homomorphisms means 
$$\varphi \iota(\alpha) = \sum_{a \in A}r_a \varphi(a) = \sum_{\varphi(a) \in \varphi(A)} r_{\varphi(a)} \varphi(a).$$
The last equality makes sense even when $\varphi \iota$ is not an injection on $A \subset M$; 
$r_{\varphi(a)}$ can be identified with the sum of all $r_{a'}$ for all $a' \in A$ such that $\varphi(a') = \varphi(a)$. 
The equality above shows that $\varphi\iota(\dirsum{R}{A})$ along with the natural injection
$$j: \varphi(A) \hookrightarrow \varphi\iota(\dirsum{R}{A}); \quad \varphi(a) \mapsto \varphi(a)$$
satisfies the universal property for the free module on $\varphi(A)$. 
Namely, the map $J: \varphi\iota(\dirsum{R}{A}) \rightarrow \im \varphi$ induced by the inclusion map $\varphi(A) \hookrightarrow \im \varphi$ is just 
$$\sum_{\varphi(a) \in \varphi(A)} r_{\varphi(a)} \varphi(a) \mapsto \sum_{\varphi(a) \in \varphi(A)} r_{\varphi(a)} \varphi(a),$$
and this map must be surjective because $\varphi \iota: \dirsum{R}{A} \twoheadrightarrow \im \varphi$ is a surjection. 

\subsection{} 
Denote $R =  \Z[x_1, x_2, \ldots]$ and $I \defeq \<x_1, x_2, \ldots\> \lhd R$.  
Suppose $I$ is finitely generated by $f_1, \ldots, f_n \in I$. 
Even though $f_1, \ldots, f_n$ are polynomials in infintely many indeterminants, by definition of a polynomial, they are expressed in terms of finitely many indeterminants, say from $x_1, \ldots, x_N$. 
Of course, by definition of $I$, none of $f_1, \ldots, f_n$ have a constant term. 
$f_1, \ldots f_n$ generate $I$, so there are $\alpha_1, \ldots, \alpha_n \in R$ such that 
\begin{equation} \label{eq:6_14_contradiction}
  x_{N + 1} = \alpha_1 f_1 + \ldots \alpha_n f_n.
\end{equation}
However, consider the substitution map $\Phi: R \mapsto \Z[x_{N + 1}, x_{N + 2}, \ldots]$ that sends $x_1, \ldots x_N$ to 0 and is constant on $\Z$ and all indeterminants $x_{N + 1}, x_{N + 2}, \ldots$. 
$\Phi$ sends the left side of \eqref{eq:6_14_contradiction} to $x_{N + 1}$ but sends the right side of \eqref{eq:6_14_contradiction} to 0, which is a contradiction. 

\setcounter{subsection}{17} 
\subsection{} 
I use bar notation for elements in the quotient $M / N$. 
Let $M / N = \<\overline{m}_1, \ldots, \overline{m}_k\>$ and $N = \<n_1, \ldots, n_t\>$. 
The following shows that $M = \<m_1, \ldots, m_k, n_1, \ldots, n_t\>$. 
Take any $m \in M$.
Letting $\pi:M \rightarrow M / N$ be the canonical projection, 
$$\pi(m) = m + N = (r_1 m_1 + \ldots + r_k m_k) + N$$ 
for some $r_1, \ldots, r_k \in R$. 
Hence $m - (r_1 m_1 + \ldots r_k m_k) \in N$, which means 
$$m = r_1 m_1 + \ldots + r_k m_k + s_1 n_1 + \ldots + s_t n_t$$
for some $s_1, \ldots, s_t \in R$. 
This shows that any $m \in M$ is in the $R$-span of $\<m_1, \ldots, m_k, n_1, \ldots, n_t\>$. 


\end{document}
